\begin{center}
\begin{minipage}{0.86\textwidth}
\small
\noindent\textbf{\color{accent}Abstract.}
Large language models are increasingly used not only to answer but to
act: to send mail, edit records, browse, and run code on a person's
behalf. The models that plan these actions state falsehoods with
confidence and follow instructions embedded in what they read, and the
tools they invoke are third-party code. We present \sys{}, a
self-hosted platform for AI agents built on a single principle: the
model may propose any action, but holds no authority to perform one.
Every proposal is a structured request that passes through gates the
model cannot read or influence---deny-by-default access control with
two machine-checked invariants, a per-call executor that holds each
agent to a manifest approved before its code runs, a person-chosen
trust level with approvals bound to the exact call by hash, and a
layered sandbox (a uid per agent, Landlock, a seccomp denylist, and an
egress proxy that admits only declared hosts). Credentials are
decrypted by the backend and lent to one agent for one call, never to
the model; what is displayed as data beside the model's words is
rendered from stored, schema-validated results of real calls rather
than from the model's text. The design applies operating-system ideas
to agents: processes, users, system calls, a firewall, a scheduler,
kill, and an append-only log. We describe the threat model, each
mechanism with the code that implements it, and an evaluation through
more than two thousand automated tests---including a deliberately
hostile agent that attempts each known escape---and we state every
limit we know of, including those the design accepts on purpose.
\end{minipage}
\end{center}
